\documentclass[a4paper,11pt]{article}

\usepackage[ngerman]{babel}
\usepackage[top=3cm,bottom=3cm,left=3cm,right=3cm]{geometry} %NICHT IN KEMIE2


%\usepackage{microtype} %schöneres typesetting  mit [stretch=10] für nur 1% anstelle von 2%
%\usepackage{lmodern} %Latin Modern FONT
%\usepackage{pifont} %schönere Font



\usepackage{amsmath}
\usepackage{nicefrac}
\usepackage{amssymb}
%\usepackage{mathtools}

\usepackage{titlesec} %Um Section, etc bearbeiten zu können
\usepackage{setspace}
\usepackage{enumitem} %enumerate
\usepackage{multicol}
\usepackage[many]{tcolorbox}
\usepackage{hyperref}
\usepackage{arydshln}

%\renewcommand{\ttdefault}{lmtt} %Schriftart wird geändert zu Latin Modern Typewriter


\hypersetup{hidelinks,
			colorlinks=true,
			linkcolor=black,
			citecolor=miudiutex,
			urlcolor=miugray2
			}

\setlength{\parindent}{0pt}
\setcounter{section}{0}


\titlespacing*{\section}{0pt}{20mm}{6mm}
\titlespacing*{\subsection}{0pt}{6mm}{3mm}
\titlespacing*{\subsubsection}{0pt}{3mm}{0mm}



%\definecolor{miudiutex}{HTML}{2f3d39}
\definecolor{miudiutex}{HTML}{1d2926}
\definecolor{miugray}{HTML}{b4cccf}
\definecolor{miugray2}{HTML}{4d7365}
\definecolor{red}{HTML}{cf1f5c}
\definecolor{green}{HTML}{00A36C}
\definecolor{delphinidin}{HTML}{315794}
\definecolor{pelargonidin}{HTML}{b33807}
\definecolor{cyanidin}{HTML}{ba0746}
\definecolor{petunidin}{HTML}{b56fed}



\raggedcolumns %wichtig für Multicolumns


%================================
% SYMBOLS
%================================

\newcommand{\RA}{$\rightarrow$~}
\newcommand{\RL}{$\rightleftharpoons$~}
\newcommand{\HARP}{\ding{226}~}
%$\xrightarrow{2}$ %Pfeil mit 2 drüber


%================================
% SHORT COMMANDS (BOXES ; ETC)
%================================


\newcommand{\notiz}[1]{\begin{notizz}{}{}\begin{center}
			{\color{miugray2!50!miudiutex}#1} \end{center}\end{notizz}}
\newcommand{\zit}[2]{\begin{zitat*}{#1}{}\small #2 \end{zitat*}}

\newcommand{\frage}[2]{\begin{qst*}{#1}{}\small #2 \end{qst*}}

\newcommand{\unten}{\end{center} \tcblower \begin{center}}


\newcommand*\circled[1]{\tikz[baseline=(char.base)]{
            \node[shape=circle,draw,inner sep=0.5pt,miudiutex] (char) {\tiny #1};}}
            
\newcommand{\miusquare}{{\color{miugray2!50}\scriptsize $\blacksquare$}~}
\newcommand{\miusquarp}{{\color{petunidin!50}\scriptsize $\blacktriangleright$}~}
\newcommand{\niu}{\item[\miusquare]}
\newcommand{\nip}{\item[\miusquarp]}

%%%%% ABSAETZE

\newcommand{\pps}{\par \vspace*{1mm}}
\newcommand{\pp}{\par \vspace*{3mm}}
\newcommand{\ppbig}{\par \vspace*{8mm}}
\newcommand{\pphuge}{\par \vspace*{10mm}}

%================================
% FRAGE BOX
%================================

\tcbuselibrary{theorems,skins,hooks}
\newtcbtheorem{qst}{Frage:}
{%
	enhanced,
	breakable,
	colback = gray!15!white,
	frame hidden,
	boxrule = 0sp,
	borderline west = {2.5pt}{0pt}{red!70},
	sharp corners,
	detach title,
	before upper = \tcbtitle\par\smallskip,
	coltitle = black,
	fonttitle = \bfseries\sffamily,
	description font = \mdseries,
	separator sign none,
	segmentation style={solid, gray!50!black},
}
{th}




%================================
% ZITAT BOX
%================================

\tcbuselibrary{theorems,skins,hooks}
\newtcbtheorem{zitat}{Zitat}
{%
	enhanced,
	breakable,
	colback = gray!15!white,
	frame hidden,
	boxrule = 0sp,
	borderline west = {2.5pt}{0pt}{black!70},
	sharp corners,
	detach title,
%	before upper = \tcbtitle\par\smallskip,
	coltitle = gray!50!black,
	fonttitle = \bfseries\sffamily,
	description font = \mdseries,
%	separator sign none,
	segmentation style={solid, gray!50!black},
}
{th}



%================================
% SIMPLE SCHWARZER RAHMEN BOX
%================================


\newcommand{\boxfr}[1]{
	\begin{tcolorbox}[
	drop shadow southeast,
	enhanced,
	colframe=black!50,
	colback=white,
	boxrule = 1pt, 
%	toprule=0pt, bottomrule=0pt,rightrule=0pt,leftrule=0pt,
	boxsep=5pt,
	arc=1pt,
	]
	\begin{center}
	#1
	\end{center}
	\end{tcolorbox}
}

\definecolor{miudiutex}{HTML}{1d2926}
\definecolor{miugray}{HTML}{b4cccf}
\definecolor{miugray2}{HTML}{4d7365}
\definecolor{white2}{HTML}{f5f7f8}
\definecolor{red}{HTML}{cf1f5c}
\definecolor{green}{HTML}{00A36C}
\definecolor{delphinidin}{HTML}{315794}
\definecolor{uniblau}{HTML}{004e9f}
\definecolor{unigelb}{HTML}{fcba00}
\definecolor{unigrau}{HTML}{909085}
\definecolor{pelargonidin}{HTML}{b33807}
\definecolor{cyanidin}{HTML}{ba0746}
\definecolor{paeonidin}{HTML}{d15ca6}
\definecolor{petunidin}{HTML}{b56fed}
\definecolor{malvidin}{HTML}{8a2d8c}

\newcommand{\metermilli}{\mskip3mu \text{mm}}
\newcommand{\cm}{\mskip3mu \text{cm}}
\newcommand{\m}{ \text{m}}
\newcommand{\lite}{ \text{l}}
\newcommand{\kg}{ \text{kg}}
\newcommand{\g}{ \text{g}}
\newcommand{\km}{ \text{km}}
\newcommand{\h}{ \text{h}}
\newcommand{\s}{ \text{s}}
\newcommand{\tonne}{ \text{t}}
\usepackage[utf8]{inputenc}
\usepackage[T1]{fontenc}
\usepackage{lmodern}

\usepackage[
    activate={true,nocompatibility},
    final,
    tracking=true,
    kerning=true,
    spacing=true,
    factor=1100,
    stretch=30,
    shrink=20
]{microtype}
\usepackage{pdfpages}
\usepackage{awesomebox}
\usepackage{marginnote}
\tcbuselibrary{documentation}
\usepackage{sidenotes}
\usepackage{import}
\usepackage{xifthen}
\usepackage{transparent}
\usepackage[table]{xcolor}


\newcommand{\abbicon}{%
  \protect\tikz[baseline=0.2mm,scale=0.8]{%
    % Das blaue Quadrat mit abgerundeten Ecken
    \fill[uniblau, rounded corners=1.2pt] (0,0) rectangle (0.8em, 0.8em);
    % Der weiße Kreis in der Mitte
    \fill[white] (0.4em, 0.4em) circle (0.12em);
  }%
} 
\usepackage[labelfont={bf},font={color=miudiutex,small},figurename={\abbicon $~$Abbildung}]{caption}


\newcommand{\bckgrnd}{
\AddToHookNext{shipout/background}{
\begin{tikzpicture}[remember picture, overlay]
 \draw[draw=none,fill=uniblau!70, shift={(current page.north west)}] (-3,0) rectangle (23,-3.75);
 \draw[draw=none,fill=miugray, shift={(current page.north east)}] (-7,0) rectangle (3,-3.75);
\end{tikzpicture}
}
}



\newcommand{\incfig}[2]{%
\def\svgwidth{#2\columnwidth}
\import{C://LaTeX-Files/pngs/}{#1.pdf_tex}
}

\renewcommand{\textbf}[1]{\textsf{{\bfseries {#1}}}}



\titleformat{\section}[hang]
		{\LARGE \color{uniblau!75!black} \sffamily \bfseries}
		{\thesection}
		{6mm}
		{}
		[]
\titleformat{\subsection}[hang]
		{\large \sffamily \bfseries \color{black}}
		{\thesubsection}
		{4mm}
		{{\color{miugray}$\blacksquare ~$}}

\titleformat{\subsubsection}[block]
		{\normalsize \bfseries \sffamily \color{miudiutex}}%
		{\normalsize \color{black} \thesubsubsection}
		{2mm}
		{{\color{miugray}$\blacksquare ~$}}
		[ \color{black}]
		
\titleformat{\paragraph}
		{\normalsize \bfseries \sffamily \color{black}}
		{\footnotesize \theparagraph}
		{1mm}
		{}

\pagestyle{empty}
\titlespacing*{\section}{0pt}{20mm}{12mm}
\titlespacing{\section}{0pt}{20mm}{12mm}

\titlespacing*{\subsection}{0pt}{14mm}{4mm}
\titlespacing{\subsection}{0pt}{14mm}{4mm}

\titlespacing*{\subsubsection}{0pt}{6mm}{3mm}
\titlespacing{\subsubsection}{0pt}{6mm}{3mm}

\titlespacing*{\paragraph}{0pt}{6mm}{2mm}
\titlespacing{\paragraph}{0pt}{6mm}{2mm}


\let\oldmarginfigure\marginfigure
\let\endoldmarginfigure\endmarginfigure

\let\oldmarginfigure\marginfigure
\let\endoldmarginfigure\endmarginfigure

\renewenvironment{marginfigure}[1][0pt] % [1] steht für 1 Argument, [0pt] ist der Standardwert
  {\oldmarginfigure[#1]\captionsetup{labelfont={sf,bf,color=uniblau!75!black},font={color=miudiutex,small}}}
  {\endoldmarginfigure}
  
  
  
  
\renewcommand{\labelitemi}{\color{uniblau!70}$\bullet$}
\renewcommand{\labelitemii}{\color{uniblau!70}$\cdot$}
\newcommand{\plmtsquare}{{\color{uniblau!70}$\blacksquare~$}}
\newcommand{\splmt}[1]{{\color{uniblau!75!black} \sffamily #1}}
\newcommand{\fplmt}[1]{{\color{uniblau!75!black} \sffamily \bfseries #1}}

\newcommand{\antwort}[1]{
\begin{tcolorbox}[
	enhanced,
	enforce breakable,
	standard jigsaw,
	boxrule=0.7pt,
	colframe=black,
	sharp corners,
	boxsep=5pt,
	title=\bfseries \sffamily Antwort,
	opacityback=0,
	coltitle=miudiutex,
	attach boxed title to top text left={yshift=-\tcboxedtitleheight/2},
	boxed title style={colback=white,colframe=white},
	bottomrule at break=0pt,
	toprule at break=0pt,
	pad at break=16mm,
]
	\begin{center}
	#1
	\end{center}
	\end{tcolorbox}
}




\rowcolors{2}{miugray!50}{miugray!50}
\arrayrulecolor{white} % Weiße Gitterlinien
\setlength{\arrayrulewidth}{1.5pt} % Etwas dickere Linien
\renewcommand{\arraystretch}{1.4} % Mehr Zeilenhöhe für bessere Lesbarkeit}


%\setlist[itemize]{itemsep=0mm}
\setlist[enumerate]{itemsep=0mm}
\usetikzlibrary{patterns, arrows.meta, calc, matrix}

\newif\ifshowme
\showmetrue




\begin{document}
\newgeometry{top=1.8cm,bottom=4cm,left=1.5cm,right=7.5cm,marginparsep=-2.00cm,marginparwidth=4.75cm}


\bckgrnd


\begin{center}
\LARGE
\color{white}
\textbf{Übung 08} \par \vspace*{2mm}
\large \color{miugray} \textbf{Prozessbezogene Lebensmitteltechnologie}
\end{center}
\par 
\vspace*{6mm}


\color{miudiutex}
\section*{Mischungsrechnungen}
 \marginnote[]{\small \fplmt{Anmerkung:} Da die Themen in der Vorlesung nicht besprochen worden, recherchieren Sie bitte eigenständig die Gleichungen und Therme dieser Übung.}
 
%\subsection*{Aufgabe 1}
%Der pK$_S$-Wert von Buttersäure beträgt 4,82.
%\marginnote[]{\small \fplmt{Tipp:} Versuchen Sie das Massenwirkungsgesetz nach dem Anteil der undissozierten Säure umzustellen}
% Berechnen Sie den Anteil an undissoziierter Säure bei den pH-Werten 1, 2, 3, 4, 5 und 6 und 7.
%
%
%\ifshowme
%\antwort{
%
%\begin{align*}
%K_S &= \frac{[H_3O^+] \cdot [A^-]}{[HA]}  \\
%pK_S &= -\log K_S \\
%pH &= -\log [H_3O^+]
%\end{align*}
%Henderson-Hasselbalch:
%\begin{align*}
%pK_S &= pH -\log(\frac{[A^-]}{HA}) \\
% &= -\log [H_3O^+] -\log(\frac{[A^-]}{HA})
%\end{align*}
%Umstellen:
%\begin{align*}
%pH - pK_S &= \log(\frac{[A^-]}{HA}) & | 10^x \\
%10^{pH - pK_S} &= \frac{[A^-]}{HA} \\
%	&~\hat{=}~ \frac{1 - [HA]}{HA} \\
%	&~\hat{=}~ \frac{1}{HA} - \frac{[HA]}{HA} \\
%	&~\hat{=}~ \frac{1}{HA} - 1 \\
%\\
%10^{pH - pK_S}	&~\hat{=}~ \frac{1}{HA} - 1 \\
%10^{pH - pK_S} +1 &= \frac{1}{HA} \\ \\
%HA &= \frac{1}{10^{pH - pK_S} + 1 }
%\end{align*}
%Nachdem Einsetzen von pK$_S$ und pH-Wert erhält man folgende Ergebnisse: \pp
%pH 1 \RA 99,985\% \par pH 2 \RA 99,85\% \par pH 3 \RA 98,5\% \par pH 4 \RA 86,9\% \par pH 5 \RA 39,8\% \par pH 6 \RA 6,2\% \par pH 7 \RA 0,656\%
%}
%\else
%\fi


\subsection*{Aufgabe 1}
Eine Probe Cornflakes hat einen gravimetrischen Wassergehalt von 7,5\% bezogen auf die
Feuchtmasse. Berechnen Sie den gravimetrischen Wassergehalt bezogen auf die
Trockenmasse.

\ifshowme
\antwort{
\begin{align*}
\frac{7,5}{92,5}  \qquad = \qquad 0,081 \qquad \hat{=} \qquad 8,11 \%
\end{align*}

}
\else
\fi


\subsection*{Aufgabe 2}
Eine Charge Cornflakes hat einen gravimetrischen Wassergehalt von 21,2\% bezogen auf die
Trockenmasse. Berechnen Sie den gravimetrischen Wassergehalt bezogen auf die
Feuchtmasse.

\ifshowme
\antwort{
\begin{align*}
\frac{21,2}{121,2} \qquad = \qquad 17,5 \%
\end{align*}
}
\pagebreak
\else
\fi



\subsection*{Aufgabe 3}
1.200 g einer 24\%igen Lösung sollen durch Eindampfen auf 60\% konzentriert werden. Wie viel
Gramm Wasser sind zu entziehen? Nutzen Sie folgende Wege zur Berechnung:


\begin{enumerate}[label=\alph*)]
\item Mischungsgleichung
\ifshowme

\antwort{
\marginnote{\small \fplmt{Anmerkung:} $x_1 = 0\%$ da die Wassermasse keinen Mengenanteil aus der Lösung besitzt. Wasser = 100$\%$ deswegen gelöster Stoff = 0$\%$}

Mit: $x_1 = 0\%$ ; $x_2 = 60\%$ ; $x_3 = 24\%$ ; $m_1 = ?$ ; $m_2 = m_3 - m_1$ ; $m_3 = 1200$
\begin{align*}
m_1x_1 + m_3x_2 - m_1x_2 &= m_3x_3 \\
m_1x_1 - m_1x_2 &= m_3x_3 - m_3x_2 \\
m_1(x_1-x_2) &= m_3(x_3-x_2) \\
m_1 &= \frac{m_3(x_3-x_2)}{(x_1-x_2)} \\
m_1 &= \frac{1200(0,24-0,6)}{0-0,6}\\
 &= 720
\end{align*}
Daher müssen 720 g Wasser entzogen werden. 
}
\else
\fi

\item Mischungskreuz 
\ifshowme
\antwort{
%
%(Beispiel aus 35\%iger Lösung und 22\%iger Lösung, NICHT LÖSUNG DER AUFGABE) \par 
%\begin{tabular}{lllll}
% &Gehalt&Teile&Gehalt&Teile \\ \cline{1-5}
% Lösung &x&z-y&35\% &22 - 0 = 22 \\
% Mischung & z &&22\% & \\
% Lösung &y& x-z & 0\% & 35 - 22 = 13
%\end{tabular}

\begin{tikzpicture}
\matrix (m) [
    matrix of math nodes, 
    row sep=1.5em, 
    column sep=2.5em, 
    ampersand replacement=\&
  ] {
    60\% \&      \& \vert 24 - 0 \vert = 24 \\
         \& 24\% \& \\
    0\%  \&      \& \vert 60 - 24 \vert = 36 \\
  };
  % Pfeile von links zur Mitte
  \draw[->, thick] (m-1-1) -- (m-2-2);
  \draw[->, thick] (m-3-1) -- (m-2-2);

  % Pfeile von der Mitte nach rechts
  \draw[->, thick] (m-2-2) -- (m-1-3);
  \draw[->, thick] (m-2-2) -- (m-3-3);
\end{tikzpicture}
\par \vspace*{-1cm}
\begin{align*}
60~\text{Teile} ~&\hat{=}~ 1200 ~\text{g} \\
1~\text{Teil} ~&\hat{=}~ 20 ~\text{g} \\
36~\text{Teile} ~&\hat{=}~ 720 ~\text{g} \\
\end{align*}
}
\else
\fi

\item Prozentrechnung
\ifshowme
\antwort{
\begin{align*}
1200~\text{g} \cdot \frac{24}{60} &= 480~\text{g} \\
1200~\text{g} \cdot \frac{36}{60} &= 720 ~\text{g}
\end{align*}
}
\else
\fi

\end{enumerate}

%\pagebreak
\subsection*{Aufgabe 4}
Wie viel kg 10\%ige Saccharose-Lösung müssen zu 100 kg 30 \%iger Saccharose-Lösung
gegeben werden, um eine 20\%ige Saccharose-Lösung zu erhalten? Berechnen Sie die Lösung
über die Mischungsgleichung.

\ifshowme
\antwort{

Mit: $x_1 = 10\% $ ; $x_2 = 20\%$ ; $x_3 = 30\% $ ; $m_1 = ? $ ; $m_2 = m_3 + m_1$ ; $m_3 = 100$ 
\begin{align*}
m_3x_3 + m_1x_1 			&= m_3x_2 + m_1x_2 \\
m_3x_3 + m_1x_1 - m_1x_2 	&= m_3x_2 \\
m_1x_1 - m_1x_2 			&= m_3x_2 - m_3x_3 \\
m_1(x_1 - x_2) 				&= m_3(x_2-x_3) \\
m_1 						&= \frac{m_3(x_2-x_3)}{x_1 - x_2} \\
m_1 						&= \frac{100(0,2-0,3)}{0,1 - 0,2} \\
&= 100
\end{align*}
}
\else
\fi

%\subsection*{Aufgabe 5}
%\tcbdocmarginnote[colback=unigelb!50,colframe=uniblau]{
%Wiederholungsaufgabe
%}
%Sie haben den Erhitzungsprozess so entwickelt, dass bei einer Temperatur von 121 °C eine
%Reduktion der Keimzahl um 99,999\% erfolgt. Der z-Wert der Mikroorganismenabtötung
%beträgt 8 °C. Aufgrund eines Fehlers in der Regelung der Temperatur liegt die reale
%Temperatur bei 120 °C. 
%\begin{itemize}
%
%\item Berechnen Sie, um wieviel Prozent die Keimzahl tatsächlich reduziert
%wird.
%\ifshowme
%\antwort{
%
%\begin{align*}
%D_1 &= D_2 \cdot 10^{\frac{T_1-T_2}{z}} \\
%D_1 &= 5 D \cdot 10^{\frac{120-121}{8}} \\
%D_1 &= 3,75 \\
%\\
%1 - 1 \cdot 10^{-3,75} &= 0,99982
%\end{align*}
%}
%\else
%\fi
%
%\item Berechnen Sie den Effekt, wenn ein 2D-Konzept umgesetzt wird.
%\ifshowme
%\antwort{
%\begin{align*}
%D_1 &= D_2 \cdot 10^{\frac{T_1-T_2}{z}} \\
%D_1 &= 2 D \cdot 10^{\frac{120-121}{8}} \\
%D_1 &= 1,4998 \\
%\\
%1 - 1 \cdot 10^{-1,4998} &= 0,96836
%\end{align*}
%
%
%}
%\else
%\fi
%
%\end{itemize}


\subsection*{Aufgabe 5}
Für eine Mikrowellenfrequenz von 2450 MHz und eine Temperatur von \mbox{20 °C} beträgt die
stoffabhängige Permitivitätszahl $\varepsilon '$ von rohen Kartoffeln 64 und der Verlustwinkel tan $\delta$ 0,23.
Nehmen Sie als Lichtgeschwindigkeit den gerundeten Wert 3 * 10$^8$ m/s an.
\begin{enumerate}[label=\alph*)]

\item Bestimmen Sie anhand untenstehender Formel die Eindringtiefe des Feldes in die
Kartoffel.
\item Bestimmen Sie die Eindringtiefe auch für Wasser bei 25°C ($\varepsilon '$ 78; tan $\delta$: 0,16) und Eis
($\varepsilon '$: 3,2; tan $\delta$: 0,0009).
\item Wie ist die Eindringtiefe definiert?

\end{enumerate}
Es gilt für die Eindringtiefe Z mit der Wellenlänge $\lambda$: 
\begin{align*}
Z = \frac{\lambda}{2 \pi} \left[ \frac{2}{\varepsilon ' \cdot (\sqrt{1+ \tan^2 \delta - 1)}} \right]^{\frac{1}{2}}
\end{align*}

\ppbig 


\ifshowme
\antwort{
Für das Lösen der oben stehenden Aufgaben müssen wir zunächst die Wellenlänge $\lambda$ ausrechnen:


\begin{align*}
\lambda = \frac{c}{f} \qquad \rightarrow \qquad \frac{3 \cdot 10^8 ~\text{m/s}}{2.450 \cdot 10^6 ~\text{Hz}} = 0,1224 ~\text{m}
\end{align*}

\marginnote[]{\small \fplmt{Anmerkung:} Vorsicht beim Einsetzen von tan $\delta$. Die Darstellung $\tan^2$ bezieht sich auf das Quadrat des Tangens und nicht von $\delta$}[1cm]
Dann Werte einsetzen (Kartoffel): 

\begin{align*}
Z &= \frac{0,1224 ~\text{m}}{2 \pi} \left[ \frac{2}{64 \cdot (\sqrt{1+ (0,23)^2 - 1)}} \right]^{\frac{1}{2}} \\
 &= 0,0213 ~\text{m}
\end{align*}

Das Selbe für Wasser: 

\begin{align*}
Z &= \frac{0,1224 ~\text{m}}{2 \pi} \left[ \frac{2}{78 \cdot (\sqrt{1+ (0,16)^2 - 1)}} \right]^{\frac{1}{2}} \\
 &= 0,0276 ~\text{m}
\end{align*}


Und Eis:

\begin{align*}
Z &= \frac{0,1224 ~\text{m}}{2 \pi} \left[ \frac{2}{3,2 \cdot (\sqrt{1+ (0,0009)^2 - 1)}} \right]^{\frac{1}{2}} \\
 &= 24,2 ~\text{m}
\end{align*}


}
\else
\fi




\ifshowme
\antwort{
\textbf{Eindringtiefe Z:} \par Punkt, an dem die elektrische Feldstärke nur noch 1/e von der ursprünglichen Feldstärke entspricht.

}
\else
\fi














\end{document}